%% ============================================================================
%%  sobre1-comun.tex -- Datos y utilidades comunes del Sobre N.° 1
%%
%%  Lo cargan todos los documentos del sobre después de \documentclass.
%%  Los datos de la empresa (razón social, RUT, domicilio, representante,
%%  cargo, correo y teléfono) vienen de lafrox.cls, sección 6, iguales a los
%%  de todas las portadas de la propuesta.
%% ============================================================================

\datoslafrox{
  caso       = {Caso --- Logística},
  fecha      = {25 de noviembre de 2026},
  documento  = {Antecedentes del Oferente},
  sobre      = {Sobre N.\textdegree{} 1 \textemdash{} Antecedentes del Oferente},
}

% El Sobre N.° 1 se entrega físico y foliado: el número se rotula como folio.
\renewcommand{\LFXfolio}{Folio}

% ---- Antecedentes del Formulario A-1 que no están en lafrox.cls -------------
\newcommand{\LFXnombrefantasia}{LafroX}
\newcommand{\LFXgiro}{Desarrollo, integración y operación de sistemas para logística, distribución de consumo masivo y transporte de carga}
\newcommand{\LFXsitioweb}{\url{https://lafrox.github.io/}}
\newcommand{\LFXciudad}{Valparaíso, Región de Valparaíso}
\newcommand{\LFXindustria}{empresa distribuidora de consumo masivo, mandante Distribuidora Puelche S.A.}
\newcommand{\LFXcedularepresentante}{21.529.759-4}
\newcommand{\LFXcontraparte}{Alex Aravena, Jefe de Proyecto}
\newcommand{\LFXcorreocontraparte}{a.aravena@lafrox.cl}
\newcommand{\LFXtelefonocontraparte}{+56 32 250 4100}

% ---- Foliación correlativa del sobre ---------------------------------------
%  folios.tex lo escribe compilar_sobre1.py: define el folio de inicio y de
%  término de cada documento y las filas del índice del Formulario A-5.
\makeatletter
\newcommand{\definirfolio}[3]{%
  \expandafter\gdef\csname LFXfi@#1\endcsname{#2}%
  \expandafter\gdef\csname LFXft@#1\endcsname{#3}}
\newcommand{\filasindicesobre}{}
\InputIfFileExists{folios.tex}{}{}
\edef\LFXfolioinicial{\ifcsname LFXfi@\jobname\endcsname
  \csname LFXfi@\jobname\endcsname\else 1\fi}

% ---- Tabla cuyas filas vienen en una macro ---------------------------------
%  tablalafrox lee su cuerpo como argumento y le agrega los filetes de fila;
%  para que vea filas reales, la macro de filas se expande una vez antes.
%  \tablaconcuerpo{caption}{label}{columnas}{cabecera}{\macrodefilas ...}
\newcommand{\tablaconcuerpo}[5]{%
  \expandafter\LFX@tablaconcuerpo\expandafter{#5}{#1}{#2}{#3}{#4}}
\newcommand{\LFX@tablaconcuerpo}[5]{%
  \begin{tablalafrox}{#2}{#3}{#4}{#5}#1\end{tablalafrox}}
\makeatother

% ---- Apertura del formulario -----------------------------------------------
%  \aperturaformulario{A-1}{Identificación del Proponente}
\newcommand{\aperturaformulario}[2]{%
  \renewcommand{\thetable}{#1.\arabic{table}}%
  \chapter*{Formulario #1: #2}%
  \addcontentsline{toc}{chapter}{Formulario #1: #2}%
  \markboth{Formulario #1}{}}

% ---- Firma del representante legal -----------------------------------------
%  Nombre, firma y fecha, como cierran los formularios A-2 a A-6.
\newcommand{\bloquefirma}{%
  \par\addvspace{20pt}%
  \noindent\begin{minipage}{\textwidth}
    \begin{tabular}{@{}l@{\hspace{8pt}}l@{}}
      Nombre: & {\intersemibold\LFXrepresentante}, \LFXcargo, \LFXempresa\\[10pt]
      Firma: & \raisebox{-0.45\height}{\includegraphics[width=3.4cm]{\LFXfirmaimagen}}\\[6pt]
      Fecha: & \LFXfecha
    \end{tabular}
  \end{minipage}\par}
